% v12 Background and Related Work (condensed from v11 Section II).
\section{Background and Related Work}\label{sec:related}

\textbf{Robust annotation anchoring.} Work on annotation positioning after document edits \cite{brush2001robust}, robust intra-document locations \cite{phelps2000robust}, quote-based fuzzy anchoring \cite{csillag2013fuzzy} and the W3C text-quote and text-position selectors \cite{w3c2017annotation} anchors annotations by content. For code, Meta-Manager records edit provenance in the editor \cite{horvath2024metamanager}, Magic Markup re-tags markup with a language model \cite{misback2024magic}, and Codetations combines tracked edits with model resolution of offline changes \cite{misback2025codetations}. These follow one developer's working copy and its edit stream; \sys{} receives whole submitted versions and must decide, for each earlier concern, whether its supervisor needs to look again.

\textbf{Code review and warning tracking.} Modern review combines line-specific discussion with human judgment \cite{bacchelli2013expectations,sadowski2018modern}. GitHub marks a comment outdated when its diff hunk changes \cite{github2018outdating}; Gerrit ports unresolved comments to a calculated position, or to file level, and ported comments look like ordinary ones \cite{gerritporting}; GitLab can resolve outdated threads automatically \cite{gitlabautoresolve}, the opposite of human-only resolution. CodeReviewQA tests whether language models localize the change a comment asks for within one version \cite{lin2025codereviewqa}; \sys{} must find where reviewed code went in the next. Static-analysis trackers keep warning identity and review decisions across versions \cite{avgustinov2015tracking,li2024tracking,spacco2006tracking,sonarqube}, and ViolationTracker builds precise violation histories with introduced, fixed, deleted and reopened states \cite{yu2023violationtracker}. Its lifecycle is the closest model of a persistent issue, but a violation is \emph{fixed} when the analyzer stops reporting it; a supervisor's concern has no analyzer, so \sys{} never infers closure from a disappearance.

\textbf{Change impact versus invalidation.} Anand separates two questions for stored claims about a repository: whether a diff touches what a claim depends on, and whether the claim is now false; content anchoring answers the first, fires on most changes and is a poor proxy for the second \cite{anand2026impact}. A \sys{} change label is of the first kind. We therefore use it only to ask the supervisor to look again and never to close or keep a concern, and Section~\ref{sec:suggest} tests a claim-relative question as an optional, non-resolving signal.

\textbf{Differencing and region mapping.} Line differencing underlies most revision tools \cite{myers1986ond}; our matcher uses Python's sequence matcher, an extension of Ratcliff--Obershelp \cite{ratcliff1988pattern,pythonsequence}. LHDiff tracks source lines by content and context \cite{asaduzzaman2013lhdiff}, other work tracks changed lines and program elements \cite{canfora2007identifying,kim2006program,reiss2008tracking}, ChangeDistiller and GumTree difference syntax trees \cite{fluri2007change,falleri2014fine}, refactoring-aware differencing adds multi-mappings \cite{alikhanifard2024astdiff}, and origin analysis handles renamed, merged and split entities \cite{godfrey2005using}. History trackers follow named methods, variables and blocks \cite{grund2021codeshovel,jodavi2022codetracker,hasan2024blocktracking,islam2025historyfinder}. CodeTracker's block tracking gives a block a structural identity through refactorings \cite{hasan2024blocktracking}, whereas a review anchor is an arbitrary line range with no identity of its own; it could serve as a locator under the layer for anchors that coincide with a block, but it is implemented for Java and is not run here. ReCite repairs kernel comments whose referenced function was renamed or removed \cite{sun2026recite}. Closest is CodeMapper, which maps an arbitrary region between commits by combining diff algorithms, movement detection and text search \cite{hu2026codemapper}. These methods answer \emph{where is this code now}; none defines what a reviewer should be told about the change, or what happens to a concern when the answer is uncertain.

\textbf{What is conventional and what is new.} Each element of the layer has a precedent: people resolve threads in GitHub and Gerrit, GitHub's outdating signals change conservatively, Gerrit ports positions, warning trackers keep identity and dismissals, and CodeMapper maps regions. Table~\ref{tab:related} scores the closest systems against S1--S5 (Section~\ref{sec:props}). None combines relocation with a request computed from the reviewed text, none marks a placement as ambiguous, and none treats a disappearance as a reason to ask rather than to close. The contribution is that conjunction as a locator-independent contract, and its measurement over five locators. Feedback research motivates specific, actionable comments \cite{hattie2007power,shute2008focus}, and programming-feedback tools support diagnosis more than planning corrections \cite{keuning2018systematic}; we know of no study showing that relocated annotations improve students' response to feedback.

\begin{table*}[t]
\caption{Requirements S1--S5 (Section~\ref{sec:props}) Against the Closest Systems}\label{tab:related}
\centering\footnotesize
\setlength{\tabcolsep}{3pt}
\begin{tabular}{@{}p{.17\textwidth}p{.105\textwidth}p{.105\textwidth}p{.105\textwidth}p{.105\textwidth}p{.12\textwidth}p{.105\textwidth}p{.105\textwidth}@{}}\toprule
Requirement & GitHub \cite{github2018outdating} & Gerrit \cite{gerritporting} & CodeMapper \cite{hu2026codemapper} & CodeTracker \cite{hasan2024blocktracking} & ViolationTracker \cite{yu2023violationtracker} & Claim check \cite{anand2026impact} & \sys{} \\ \midrule
Anchor is an arbitrary line range & yes & yes & yes & no: block & no: warning & no: stored claim & yes \\
Follows edited code & no: outdated & calculated position & yes & yes & yes, by matching & n/a & yes, any locator \\
S1: signals change since review & yes, by hiding the position & no: ported comments look unchanged & no review state & change history of the block & n/a & asks if the claim holds & request from $q_0$ \\
S2: flags a doubtful placement & n/a & file level if no position & no & no & no & n/a & copy guard: ambiguous \\
S3: absence is visible & outdated & file level & no match & block removed & deleted state & n/a & \lbl{removed}, request \\
S4: only a person closes & yes & yes & n/a & n/a & no: fixed when unreported & no: automatic invalidation & yes (one exception) \\
S5: concern outlives its code & thread stays, outdated & unresolved ported & n/a & n/a & no & n/a & not yet: three gaps \\
Safety outcomes measured & none found & none found & accuracy only & accuracy only & matching accuracy & claim precision & this paper \\ \bottomrule
\end{tabular}
\tablenote{}{Product cells summarize the cited documentation or paper, as accessed in October and November 2026, not measured behavior; the only product behavior measured here is GitHub's outdated flag (Section~\ref{sec:suggest}). ``None found'': no published evaluation of the product's relocation was found. The \sys{} exception to S4 is dismissal memory, which creates a matching new finding as \lbl{dismissed}; the S5 gaps are issue deletion, dismissal memory and uncomposed versions, which a proposed contract closes in a bounded model (Section~\ref{sec:statecheck}).}
\end{table*}
